% =============================================================
% §8 结论
% =============================================================
\section{结论}
\label{sec:conclusion}

本文研究了云--边 LLM 推理场景下的 AIGC 感知强化学习调度。我们构建了一个建模五项物理特性---模型权重驻留、KV cache 本地性、连续批处理、显存带宽下限、两阶段请求生命周期---的综合 AIGC 仿真器，随后评估了 11 种调度器，涵盖负载均衡启发式、经典与元启发式 DAG 调度器、PerLLM 风格受限 bandit（CS-UCB~\cite{yang2024perllm}）、多目标 NSGA-II 搜索~\cite{yu2025llmrouting}，以及三种 RL 变体（MLP-PPO、带 GRU 编码器的 A3C-R2N2~\cite{tuli2022a3cr2n}、带 Decima 风格~\cite{mao2019decima} GNN 编码器的 PPO）。

本文 AIGC 感知 PPO 调度器在（\slo{} 达成率、每 token 能耗）平面上对 10 个强基线取得\textbf{均值性能上的 Pareto 占优}：典型配置（edge${=}5$、$\lambda{=}2$）下相对最强单目标基线 \slo{} 提升 $+28\%$、能耗下降 $-7\%$，且 \slo{} 维度对全部 5 个单目标对比方法在 $p<0.05$ 水平显著。最强基线---knee 点 NSGA-II 搜索---在典型负载下与本文方法持平，在中等负载甜点区被显著超越（$p<0.01$），在每种配置下每 token 能耗都更高，并多花 $5.4\times$ 的调度时间算力。该 Pareto 占优在边缘数 $\{3, 5, 7\}$ 范围内稳健延伸，并经受住了隔离 AIGC 贡献与通用 RL 机制的严格 A3C-R2N2 受控对比。

然而 12 组件消融揭示：\textbf{AIGC 感知调度是系统级整体贡献，而非精巧的奖励技巧}---仅连续批处理仿真（消融后 $-83\%$ \slo{}）与充足预训练（$-38\%$）表现为显著组件；所有单项 AIGC 奖励与状态组件在统计上与完整模型不可区分，且在四组备选奖励权重向量下重新训练，两个目标均保持不变。这警示社区：在缺乏联合消融时，不应声称特定 AIGC 奖励项带来增益。

我们希望本工作支持两条更长期的方向：基于所开源平台构建的\textbf{真实 AIGC 调度基准}，以及面向 AIGC 调度论断的\textbf{积极联合消融文化}。
